\chapter{潜在结果、估计目标与识别}
\label{ch:potential-identification}

\begin{learninggoals}
读完本章，你应当能够：用潜在结果定义处理效应和目标总体；区分一致性、可交换性、正值性与无干扰假设；从反事实目标推导可观测表达式，并诊断选择、测量、缺失与外推造成的识别失败。
\end{learninggoals}

\section{潜在结果：用替代状态定义因果效应}

\subsection{潜在结果、基本难题与估计目标}

潜在结果框架把因果比较写进结果变量本身。对单位 $i$，$Y_i(1)$ 表示该单位接受处理时的结果，$Y_i(0)$ 表示未接受处理时的结果。个体因果效应是
\[
\tau_i=Y_i(1)-Y_i(0).
\]
这里的“处理”不必是药物，也可以是政策、教育方案、暴露水平或信息呈现；但两个处理状态必须足够明确，能够说明它们何时开始、持续多久、如何实施。

这种表述继承Neyman在随机试验中的思想，并由Rubin等推广到更一般的因果推断\parencite{neyman1923,rubin1974,holland1986}。它的核心不是某一种匹配算法，而是先定义同一单位在互斥替代状态下的结果，再讨论这些不可同时观察的量如何由研究设计识别。

\paragraph{因果推断的根本问题}

对同一单位和同一时点，我们最多观察 $Y_i(1)$ 或 $Y_i(0)$ 之一。令实际处理为 $A_i\in\{0,1\}$，观测结果满足
\[
Y_i=A_iY_i(1)+(1-A_i)Y_i(0).
\]
缺失的潜在结果不是普通漏填数据：我们不能回到同一历史起点让同一个人同时接受另一处理。因此，个体效应通常不可直接观察。这是\term{因果推断的根本问题}。

根本问题不意味着因果效应永远不可知。总体平均效应可以在随机化或其他可交换条件下识别：不同单位为彼此提供反事实信息。代价是从个体确定性问题转向人群分布问题，并依赖单位之间在处理赋值上的可比性。

\paragraph{估计目标不是只有一个}

最常见的\term{平均处理效应}为
\[
\operatorname{ATE}=E\{Y(1)-Y(0)\}.
\]
但研究问题可能要求其他目标：
\begin{align*}
\operatorname{ATT}&=E\{Y(1)-Y(0)\mid A=1\},\\
\operatorname{ATC}&=E\{Y(1)-Y(0)\mid A=0\},\\
\operatorname{CATE}(x)&=E\{Y(1)-Y(0)\mid X=x\}.
\end{align*}

ATE回答在目标人群中普遍实施与普遍不实施的平均差异；ATT回答实际接受者若不接受会怎样；CATE描述协变量为 $x$ 的亚组。三者在人群组成和效应异质性存在时不同。算法输出一个数之前，研究者必须说明目标是哪一个。

对于二元结果，常用尺度包括风险差
\[
E\{Y(1)\}-E\{Y(0)\},
\]
风险比和优势比。效应尺度会改变“无交互”和异质性的表现。风险差上的零效应不等于风险比上的一；优势比具有非可合并性，即使没有混杂，条件优势比与边缘优势比也可能不同。选择尺度应服从决策与科学问题，而非只看软件默认。

\subsection{一致性、干扰与随机化可比性}

\term{一致性}通常写为：若 $A_i=a$，则 $Y_i=Y_i(a)$。这看似只是记号，实则要求观察到的处理版本与潜在结果中定义的干预相对应。若“接受治疗”包括不同剂量、依从性和联合用药，$Y(1)$ 可能没有唯一含义。

处理定义需要包括：启动规则、具体版本、随访起点、允许的后续策略和结果测量时点。把肥胖、教育程度或种族直接当作可设定开关尤其困难，因为“将变量设为另一值”可能对应许多不同干预路径。此时应优先定义可执行或至少机制明确的干预，如提供某项教育资源，而不是含混询问“教育本身的效应”。
一致性也依赖无隐藏的多版本处理。不同版本若对结果相同，可以在问题目的下合并；若效果不同，就应扩展处理变量或把目标改写为随机策略效应。

\paragraph{SUTVA与单位间干扰}

\term{稳定单位处理值假设}（SUTVA）通常包括两个方面：处理没有未说明版本；一个单位的结果不受其他单位处理影响。后者可写为 $Y_i(a_i)$，而不是依赖全体赋值向量的 $Y_i(\mathbf a)$。

疫苗、社交网络、教育班级和市场政策经常违反无干扰。一个人接种会改变他人感染风险，平台向部分用户展示内容会改变其朋友体验，价格政策会通过均衡影响未处理者。此时“处理效应”不能只比较个体自己的 $a_i$，还要定义暴露映射，例如个体处理与邻居处理比例，并区分直接效应、溢出效应和总体效应。
SUTVA不是越强越好的技术便利。若科学问题本就关心传播或均衡效应，应显式建模干扰，而不是把它当作小偏差。单位边界也需与机制匹配：以学校而非学生为随机化单位，可能使校内干扰成为处理的一部分。

\paragraph{随机化如何提供可比性}

在理想随机试验中，处理赋值与潜在结果独立：
\[
A\indep\{Y(1),Y(0)\}.
\]
因此
\begin{align*}
E(Y\mid A=1)&=E\{Y(1)\},\\
E(Y\mid A=0)&=E\{Y(0)\},
\end{align*}
从而
\[
E(Y\mid A=1)-E(Y\mid A=0)=\operatorname{ATE}.
\]

逻辑关键不是“随机化使所有协变量在样本中完全相同”。有限样本仍可能不平衡。随机化保证赋值机制已知，并在重复随机化意义下使处理与潜在结果独立，使差异估计具有无偏性或可控制的随机化分布。
随机化也不自动解决依从性、失访、测量错误、干扰和外部效度。若部分被分配治疗者未接受治疗，按分配比较得到的是意向治疗效应；若想估计实际接受治疗的效应，需要额外假设或工具变量框架。试验内部效度与结果能否推广到其他人群是两件事。

\subsection{动态处理、异质性与连续暴露}

现实处理常随时间变化。令 $A_t$ 为时点 $t$ 的处理，$\bar A_t=(A_0,\ldots,A_t)$ 为处理历史，潜在结果写为 $Y(\bar a_T)$。时间变化协变量 $L_t$ 可能同时受先前处理影响，并影响以后处理与结果。若在普通回归中直接调整 $L_t$，既可能阻断部分处理效应，也可能引入碰撞点偏差。
纵向g公式、边际结构模型和结构嵌套模型用于这类问题\parencite{robins1986,hernanrobins2020}。在进入方法前，必须先统一\term{时间零点}：资格判定、处理分配和随访开始应当对齐。若把必须存活一段时间才能进入“处理组”的人从早期开始计入，会产生不死时间偏差。
失访可视为一种选择过程。若删失与潜在结果在给定历史后独立，可用逆概率删失权重等方法；若失访由未测因素决定，数据本身无法保证校正。报告失访原因和敏感性分析与选择估计器同样重要。

\paragraph{个体效应与异质性}

虽然 $\tau_i$ 通常不可识别，我们可以估计某些可观测亚组的CATE。但预测谁获益不等于恢复每个人确定的个体效应。即使准确估计 $E\{Y(1)-Y(0)\mid X=x\}$，同一 $X=x$ 亚组内仍可能存在未解释异质性。
数据驱动寻找亚组还会产生多重比较和过拟合。可靠异质性分析应预先区分验证性与探索性目标，使用样本分割或诚实估计，并报告不确定性。更重要的是，亚组差异可能反映基线风险而非机制改变；应在适当效应尺度上解释。

\begin{keypoint}
潜在结果不是“数据库中缺失的一列”，而是由处理定义、单位边界、时间结构和干预版本共同构造的反事实量。若这些概念没有定义清楚，再精确的估计也只是在计算一个含义不明的差值。
\end{keypoint}

\paragraph{多水平与连续处理}

处理不必二元。对剂量 $a\in\mathcal A$，每个单位有潜在剂量---反应曲线 $Y_i(a)$，目标可以是
\[
E\{Y(a)-Y(a')\}
\quad\text{或}\quad
\frac{\partial}{\partial a}E\{Y(a)\}.
\]
连续处理的正值性更严格：每个协变量层不可能在有限数据中拥有所有精确剂量，只能依赖平滑和局部支持。把剂量线性放进回归等于假定整条反应曲线形状，需用领域知识和灵活模型检查。

多版本问题在连续处理中更突出。同样名义剂量因给药速度、制剂和依从不同而作用不同。可把处理扩展为向量 $(a,v)$，或把研究问题定义为现实策略产生的版本混合。后者的效应依赖版本分布，迁移到另一系统时可能改变。

\subsection{复合处理、主层与潜在结果表}

若同时研究药物 $A$ 与心理干预 $B$，潜在结果写为 $Y(a,b)$。加法尺度上的交互可表示为
\[
E\{Y(1,1)-Y(1,0)-Y(0,1)+Y(0,0)\}.
\]
非零说明一种处理的效果随另一处理状态改变。交互是尺度相关的：加法交互存在时，乘法交互可能不存在，反之亦然。

析因随机试验能有效估计主效应与交互，但前提是处理组合可共同实施且没有隐藏版本。若药物改变参与心理干预的依从性，分配组合的意向治疗效应仍清楚，实际接受组合的效应则需要依从建模。
“协同机制”与统计交互也不相同。机制上两因素共同必要可能在某尺度表现强交互，但人群背景分布、结果阈值和测量会改变统计形态。不能仅凭回归乘积项显著就断言发现生物协同。

\paragraph{主层与处理后分层}

研究者有时想比较“无论分配哪组都会存活的人”中的生活质量。存活是处理后变量，直接只分析存活者会比较不同构成的人群。\term{主层}按联合潜在中间状态 $(S(1),S(0))$ 定义，例如始终存活者、仅处理下存活者等。
主层在概念上避免条件化实际处理后变量造成的混合，却通常不可完全观察：每人只看到一个 $S(a)$。识别需要单调性、排除限制或模型，且目标人群“始终存活者”难以在实践中定位。主层效应适合某些截断死亡、依从与替代终点问题，不是处理后偏差的免费修复\parencite{frangakisrubin2002}。

\paragraph{潜在结果表的思维实验}

设四人的潜在结果分别为
\[
(1,0),\quad (1,1),\quad (0,0),\quad (0,1),
\]
其中前项为 $Y(1)$，后项为 $Y(0)$。个体效应依次为 $1,0,0,-1$，ATE为0。若只报告ATE，会遗漏一个获益者和一个受害者；若处理赋值偏向本来会在处理下成功者，观察组差异还可能为正。
这个表说明平均零不等于人人无效，也说明效应异质性与选择偏差是不同问题。随机化可使平均比较识别ATE，却不能显示每人的两列；额外协变量可解释部分异质性，但不能一般恢复个体效应。

\section{识别：从反事实目标到可观测分布}

\subsection{识别、标准化与可交换性}

一个因果目标若能在给定假设下唯一写成观测数据分布的函数，就称为\term{可识别}。识别是在总体层面连接反事实量与可观测量的逻辑问题。\term{估计}则是在有限样本中近似这个函数，并量化随机误差。两者的次序不能颠倒。
例如，若未测混杂使两个不同因果模型产生完全相同的观测分布，却给出不同ATE，则无论样本多大、机器学习预测多准，ATE仍不可由该数据识别。相反，一个目标可能在原则上可识别，但样本太小、正值性接近违反或模型过拟合，使估计极不稳定。

\paragraph{标准化公式的推导}

设 $A$ 为二元处理，$Y(a)$ 为潜在结果，$L$ 为处理前协变量。若满足条件可交换性
\[
Y(a)\indep A\mid L,
\]
一致性以及正值性，则
\begin{align*}
E\{Y(a)\}
&=\sum_l E\{Y(a)\mid L=l\}P(L=l)\\
&=\sum_l E\{Y(a)\mid A=a,L=l\}P(L=l)\\
&=\sum_l E(Y\mid A=a,L=l)P(L=l).
\end{align*}

第二行使用条件可交换性，第三行使用一致性；正值性保证每个相关 $l$ 层中确实能观察到 $A=a$。右侧完全由观测分布组成，称为\term{标准化公式}或调整公式。ATE是 $a=1$ 与 $a=0$ 两个标准化均值之差。
这个推导展示识别分析的价值：每一步都对应一项可质疑的假设。若处理定义含混，一致性失败；若某类人必然接受处理，正值性失败；若存在未测共同原因，可交换性失败。估计器不能弥补概念或设计缺口。

\paragraph{可交换性与混杂}

\term{可交换性}表示在用于比较的层次中，处理组与对照组在潜在结果分布上可以互换。随机化在设计上提供无条件可交换性；观察研究通常只能主张给定充分协变量 $L$ 后可交换。

\term{混杂}不是“协变量与处理、结果都相关”这一纯统计定义，而是处理与结果之间存在会扭曲目标因果比较的共同原因结构。一个变量是否应调整取决于图中角色：

\begin{itemize}
  \item 处理与结果的共同原因通常需要控制；
  \item 处理后的中介若用于估计总效应，通常不应控制；
  \item 碰撞点不应仅因同时预测处理和结果而控制；
  \item 共同原因的代理变量可能减少混杂，也可能引入额外测量问题；
  \item 结果的强预测因子即使不是混杂，也可能提高精度。
\end{itemize}

因此，“把所有可用变量放进回归”不是保守策略。变量选择应以时间、领域知识和因果结构为基础，再用统计方法提高估计稳健性。

\paragraph{正值性：结构缺口与数据稀疏}

\term{正值性}要求对目标人群中每个具有正概率的协变量层，都有
\[
0<P(A=a\mid L=l)<1.
\]
结构性违反发生在某些人依法或生理上绝不可能接受某处理。例如，治疗方案只适用于某疾病阶段，却把其他阶段也纳入目标人群。解决办法不是改变估计器，而是重新定义目标人群或干预。

实践性违反发生在概率理论上非零但样本中极少。逆概率权重会因此极大，估计由少数单位支配。诊断包括查看倾向得分重叠、极端权重和有效样本量。截断权重可以降低方差，却改变偏差；应报告截断规则并做敏感性分析。
正值性与研究问题密切相关。条件越细，重叠越难；但遗漏真正混杂又破坏可交换性。这里不存在纯数据驱动的万能平衡，需要在科学充分性与支持范围之间作透明取舍。

\subsection{选择、测量与缺失机制}

研究样本通常经过进入、留存和完整测量等选择。若选择 $S$ 同时受处理或其原因与结果或其原因影响，给定 $S=1$ 会打开非因果路径。例如治疗影响存活，疾病严重度也影响存活；只分析存活者可能使治疗与严重度相关，即使原始试验已随机化。
选择偏差不同于混杂：混杂在处理之前造成不可比，选择是在纳入或条件化过程中扭曲关联。二者可同时存在，处理方法也不同。若在给定已测历史后选择独立于目标潜在结果，可用选择权重或删失权重重建目标分布；若选择依赖未测因素，必须依靠界限或敏感性分析。
病例对照研究说明“选择在结果上”并非总是错误。该设计有意按结局抽样，但在适当抽样机制和目标尺度下，优势比等量仍可估计。关键不是永远不条件化碰撞点，而是目标量、采样机制和分析公式是否配套。

\paragraph{测量并非透明窗口}

误分类与测量误差可以改变因果问题本身。若处理由回忆问卷测量，结果状态可能影响暴露报告，形成差异性误差；若结局评估者知道治疗分组，测量过程可能成为处理影响观测结果的额外路径。简单地说“非差异性误差只会把效应拉向零”并不普遍成立，尤其在多分类、非线性模型和多个误测变量中。
应区分潜在构念与观测指标。例如“抑郁”可能由量表分数操作化，政策真正改变的是症状、报告行为还是诊断机会，需要在测量模型中说明。验证子样本、重复测量、负对照和定量偏差分析可以评估误差影响，但仍需外部信息。

\paragraph{缺失数据与因果缺失}

普通缺失数据术语------完全随机缺失、随机缺失、非随机缺失------描述缺失机制与数据的关系。因果研究还要问缺失发生在处理、协变量还是结果，以及缺失机制是否受处理影响。处理后缺失可能是因果路径的一部分；把所有完整案例直接分析，相当于条件化在一个选择变量上。
多重插补、逆概率加权和联合模型各依赖不同假设。插补模型应与分析目标相容，并包括预测缺失和结果的变量。任何方法都不能从数据内部验证“给定已观测变量后缺失与未观测结果独立”；这正是需要敏感性分析的地方。

\subsection{目标试验、敏感性与部分识别}

\term{目标试验}策略要求先写出希望模拟的理想随机试验方案\parencite{hernanrobins2016}：

\begin{enumerate}
  \item 资格标准；
  \item 待比较的治疗策略；
  \item 分配程序；
  \item 随访起点与终点；
  \item 结局；
  \item 因果对比；
  \item 分析计划。
\end{enumerate}

然后说明观测数据如何模拟每一项。该方法能够暴露时间零点错位、不死时间、既往使用者偏差和处理策略含混。它不让观察研究自动等于试验，也不消除未测混杂，而是使设计选择清楚可审计。

\paragraph{敏感性分析是主分析的一部分}

关键识别假设通常不可由数据完全检验。严谨报告不应只在限制部分写一句“可能存在残余混杂”，而应评估结论对违反程度的敏感性。常见策略包括：

\begin{itemize}
  \item 指定一个未测混杂与处理、结果关联的范围，计算校正后效应；
  \item 使用负对照暴露或结局发现共享偏差结构；
  \item 对失访、误分类或处理版本进行情景分析；
  \item 报告部分识别界限，而不是强迫给出点识别；
  \item 比较多种合理图与变量定义下结果是否稳定。
\end{itemize}

敏感性参数应有领域解释，而不只是抽象数值。分析的目标不是证明结论不受任何偏差影响，而是说明需要多强的未建模因素才能改变决策。

\paragraph{部分识别：不知道也可以被形式化}

当假设不足以点识别时，可以求目标的上下界。设二元结果的随机试验存在部分结局缺失，若不对缺失机制作假设，可把处理组所有缺失者分别设为0或1，得到极端风险范围；两组范围组合成风险差界限。界限可能很宽，却准确表达数据能排除什么。以最少附加假设报告可行集合，是部分识别传统的基本做法\parencite{manski1990}。
增加可信假设可逐步收窄。例如认为失访者结局风险不低于已观察者，或风险最多高出某个幅度；每项收窄都应有实质解释。与一次性宣布“随机缺失”相比，分层界限让读者看到结论由哪条假设获得精度。
未测混杂也可用界限和偏差参数处理。若能根据外部研究限制混杂与处理、结果的最大关联，就可计算校正区间。目的不是制造一个“正确调整值”，而是展示何种偏差强度足以跨过决策阈值。

\paragraph{负对照的逻辑}

\term{负对照结局}应不受处理因果影响，却与目标结局共享某些混杂或测量过程。例如研究药物对某疾病时，选择机制相似但生物上不可能受药物影响的结局。若处理与负对照仍相关，提示残余偏差\parencite{lipsitch2010}。

\term{负对照暴露}不应影响目标结局，却共享处理的混杂结构，例如在合理时间关系下选择未来暴露。负对照不是凭“结果不显著”证明无混杂；需要明确共享偏差和排除直接作用的结构。负对照本身若有未知效应，会产生误导。

更强方法可在两个负对照和代理结构下识别某些未测混杂，但承担完备性等技术条件。基础实践中，负对照至少能把“未测混杂不可检验”转化为部分可观察的模型检查。

\paragraph{识别审计与设计改进}

高质量研究应保存一张从目标到公式的审计表。第一列写目标量及目标人群；第二列写每项假设的数学形式；第三列用领域语言解释；第四列列支持证据；第五列列潜在违反；第六列写诊断或敏感性分析。例如可交换性对应“给定病情、年龄等后，治疗选择不再包含预后信息”，证据可能是临床流程，风险是医生未记录判断，敏感性分析则为负对照和偏差参数。

审计表迫使“已调整所有混杂”变成可质疑陈述。它也便于跨团队协作：领域专家审查临床流程，统计人员审查公式，数据工程人员审查变量时间，决策者审查目标人群。

\paragraph{识别失败后的研究设计}

如果图显示未测共同原因使目标不可调整识别，不应立即换更复杂模型。可以收集共同原因或可靠代理；寻找影响处理但满足排除限制的制度工具；在关键节点开展小型随机或鼓励试验；获取中介以满足前门结构；融合另一个域的实验；或把问题改为可识别的局部目标。
识别分析由此具有建设性：失败证明告诉我们哪类新信息最有价值。一个明确的不可识别结论加下一步设计，优于一个依赖不可说明机器学习外推的精确数字。

\section{从目标合同到知识边界}

\subsection{从目标合同到逐行识别证明}

只写ATE并没有完成目标定义。一个可执行的估计目标至少包括总体、处理策略、对照策略、结局、时间范围、汇总尺度和竞争事件处理。总体决定对谁求平均；策略说明何时开始、允许哪些版本与依从路径；结局必须有测量规则；尺度决定风险差、风险比、比值比或生存时间差怎样表达决策意义。任何一项变化，都可能产生不同数值而不构成统计矛盾。

例如评价降压治疗，可以比较确诊后立即启动方案与延迟三个月，也可以比较始终维持不同血压区间的动态策略。前者接近治疗启动效果，后者还包含随访中剂量调整和依从。若结局是五年卒中风险，死亡是使卒中不再发生的竞争事件；把死亡者简单删失会改变目标。若临床决策关心每千人避免多少事件，绝对风险差通常比比值比更直接。

目标合同还要说明时间零点。资格判定、策略分配和随访开始若不对齐，会产生不死时间与选择偏差。数据分析中的每行记录都应能追溯到合同中的单位和时点。先写合同不是形式主义，而是防止分析过程中随着可得变量和显著结果悄悄更换问题。

\paragraph{识别证明应当逐行书写}

识别不是一句“控制所有混杂因素”，而是一条从反事实到观测量的等式链。以二元处理为例，目标是 $E\{Y(1)\}$。先用迭代期望写成
\[
E\{Y(1)\}=\sum_l E\{Y(1)\mid L=l\}P(L=l).
\]
若给定处理前协变量满足 $Y(1)\indep A\mid L$，可把第一项换成处理组中的反事实均值；再由一致性，在 $A=1$ 的单位中以观测结果替代 $Y(1)$，得到
\[
E\{Y(1)\}=\sum_l E(Y\mid A=1,L=l)P(L=l).
\]
正值性保证每个需要的层中存在接受处理者，使右侧确实由数据分布定义。

逐行书写的价值在于每个等号都有责任主体。全概率公式来自概率论；交换反事实与处理来自因果假设；一致性来自处理和结局定义；正值性来自目标总体与数据支持。估计器只近似最后的可观测表达式，不能倒过来证明前面的交换。若某层无人接受处理，机器学习平滑产生的预测是外推，不是恢复了非参数识别。
同一习惯适用于复杂设计。纵向g公式要在每个时点交换处理历史与未来潜在结果；迁移公式要交换试验参与和目标人群结果；前门公式要逐步标明行动与观察的交换。把推导写完整，能在建模之前发现究竟缺少测量、设计还是科学知识。

\subsection{纵向问题与不可识别边界}

许多处理不是一次性开关。药物剂量随症状调整，平台推荐随点击更新，公共政策随疫情或经济指标改变。此时既往处理影响后续协变量，后续协变量又影响下一次处理和结果。若把这些协变量当普通基线混杂因素直接回归，可能阻断既往处理的中介路径；若不调整，又留下后续处理的混杂。
设处理历史为 $\bar A_t$，随时间变化的协变量为 $\bar L_t$。动态策略 $g$ 根据已有历史指定每期处理。识别需要序贯可交换性：在每个时点，给定已经观察的历史，实际处理选择与策略下未来潜在结果独立；还需要各历史下策略选择有正概率。纵向g计算、逆概率加权的边际结构模型和结构嵌套模型，是处理这种双重角色的不同道路。

这些方法的主要困难不在公式长度，而在历史定义。就诊频率可能同时是病情信息和医疗系统选择结果；上一期处理会改变它，使其既是中介又是后续混杂。研究者必须说明测量时点、决策可用信息和失访过程。离散时间过粗会把处理与协变量顺序颠倒，过细又造成稀疏历史与极端权重。时间尺度是因果假设的一部分。

\paragraph{不可识别时仍然可以推进知识}

识别失败不等于研究结束。第一种选择是重新设计：收集关键共同原因、利用随机分配或制度阈值、改变抽样以补足重叠。第二种是改变目标：从总体效应转向有支持区域的效应，从自然中介效应转向可实施的干预型效应。第三种是部分识别，用界限表达现有信息允许的范围。第四种是敏感性分析，把不可观测假设参数化并报告结论在何处改变。

重要的是不要把这些策略包装成完全识别。限制到重叠人群会得到更可靠但不同的总体；假定单调性收紧界限会增加实质承诺；负对照发现残余偏差能反驳简单模型，却通常不能唯一校正。研究贡献可以是证明某个热门问题在现有数据下不可回答，并指出哪种新信息最有价值。
在决策场景中，宽界限也可能有用。若所有合理偏差情景下政策净收益仍为正，精确点估计并非必要；若界限跨越重要伤害阈值，就应延迟部署或采用可逆试点。把不可识别性连接到行动规则，比用一个未经说明的模型填补空白更诚实。

\subsection{目标量合同：先固定问题，再证明识别}

一份可复核的目标量合同至少需要固定六项内容。以“远程工作是否改善员工健康”为例，不能在分析时才根据数据可得性决定处理、时间和人群。

\begin{center}
\small
\begin{tabularx}{\textwidth}{>{\bfseries}p{2.1cm}X X}
\toprule
合同项 & 明确定义 & 不合格写法\\
\midrule
单位与总体 & 在2026年1月1日在职、可远程履职的全职员工 & “某公司人员”\\
处理策略 & 连续十二周每周至少三天远程，与每周不超过一天比较 & “远程或不远程”\\
时间零点 & 资格确认、策略分配与随访同日开始 & 用后来的出勤定义基线组别\\
结局与时窗 & 第十二周标准化心理健康量表及病假天数 & “健康变好”\\
因果对比 & 同一目标总体两套策略下的平均结局差 & 两个自选出勤组的观测差\\
失访与干扰 & 离职、换岗、团队协作和管理者排班如何处理 & 把没有末次量表者直接删除\\
\bottomrule
\end{tabularx}
\end{center}

目标固定后，才能问观察数据是否识别它。假设$L$包含基线健康、职位、照护责任、通勤距离与团队排班，候选识别链为：给定$L$后实际策略选择与两个潜在结果可交换；每个目标$L$层都有两种出勤策略；实际符合某策略时观测结果等于该策略下潜在结果；团队成员的远程安排不通过未定义的溢出影响个人结局。

\begin{workedexample}[识别失败时如何改目标]
若所有外地职位员工都必须全远程，而所有安全敏感职位都必须到岗，则对包含这两类职位的全体ATE存在结构性正值性失败。一种合格修改是把目标收缩到制度上两种策略都可行的职位；另一种是把处理改为可执行的排班选择权，而不是实际远程天数。两种修改都得到与原ATE不同的问题，必须在结果表和摘要中重新命名。使用光滑回归为从未存在的比较填值，不算修复识别。
\end{workedexample}

处理版本也可使一致性失败。远程工作可以伴随设备补贴、弹性时间或监控软件；如果不同版本对健康的影响不同，“每周三天远程”没有唯一潜在结果。可以把设备、时间自主和监管方式写入策略包，或者分别定义多个处理版本。这种修改是在重建因果对象，不是增加一个控制变量。

\begin{failurecheck}[识别假设只出现在讨论部分]
若结果表先给出“远程工作效应”，最后才在局限中说“可能有未测混杂”，读者无法知道数字依赖哪个等式。正确顺序是在估计之前列出目标、一致性、可交换性、正值性和干扰假设，并说明每项由设计、领域知识还是统计诊断支持。局限不是主结果之后的礼貌性附言，而是数字拥有因果含义的条件。
\end{failurecheck}

\begin{chapterreview}
潜在结果以同一单位的替代处理状态定义因果效应。由于只能观察一个潜在结果，个体效应通常不可直接观察；随机化或可交换性使人群平均效应可识别。ATE、ATT与CATE回答不同问题。一致性、处理版本、无干扰和时间零点是目标量定义的一部分，而不是估计阶段的附注。 识别是在明确假设下把反事实目标写成观测分布函数；估计是在样本中计算该函数。条件可交换性、一致性与正值性共同支持标准化。混杂、选择、测量和缺失是不同机制，不应以“多变量调整”统一处理。目标试验与敏感性分析把设计和假设置于估计器之前。
\end{chapterreview}

\begin{selfcheck}
\begin{enumerate}
  \item 为“远程办公的效应”定义单位、两个处理版本、结果和时间零点。
  \item 在什么情况下ATE与ATT会明显不同？给出带效应异质性的例子。
  \item 疫苗研究为何违反简单无干扰假设？应区分哪些效应？
  \item 随机化保证什么，又不保证什么？
  \item 逐行说明标准化公式推导分别使用了哪项假设。
  \item 举出结构性与实践性正值性违反各一例，并提出不同解决方案。
  \item 为什么控制所有处理和结果的预测变量可能增加偏差？
  \item 为一项电子健康记录研究写出目标试验的七个要素。
\end{enumerate}
\end{selfcheck}

\chapterreadings{
潜在结果的经典表述见\textcite{rubin1974}与\textcite{holland1986}；系统教材为\textcite{imbensrubin2015}。主层的原始框架见\textcite{frangakisrubin2002}。纵向观察研究、目标试验和时间相关偏差可读\textcite{hernanrobins2020}。干扰问题可从\textcite{hudgenshalloran2008}和\textcite{tchetgenvanderweele2012}入手。

识别与标准化的系统教材是\textcite{hernanrobins2020}和\textcite{imbensrubin2015}。目标试验方法见\textcite{hernanrobins2016}。部分识别的经典界限见\textcite{manski1990}，负对照方法见\textcite{lipsitch2010}。选择偏差与图形解释可结合\textcite{greenland2003}阅读；未测混杂的定量敏感性分析可参阅\textcite{vanderweele2017}。

随机化在贝叶斯因果推断中的作用见\textcite{rubin1978}。一致性不是一句符号约定：其基本声明与争议可对照\textcite{colefrangakis2009}和\textcite{vanderweele2009}，处理必须可定义以及同一标签可能包含多个版本的问题分别见\textcite{hernantaubman2008}与\textcite{vanderweelehernan2013}。选择机制的结构分析见\textcite{hernan2004selection}，时间变化处理与混杂的边际结构模型入口是\textcite{robins2000msm}。正值性在实际数据中的含义见\textcite{westreichcole2010}；如何在分析之前以设计减少观察研究偏差，可续读\textcite{rosenbaum2010}。
}
